% O1, formalización autocontenida del modelo histórico V1.
% Fuente ejecutable: specs/tla/CrossShardCommit.tla (paper1-q3-v1).
% No atribuir los 1272 trabajos históricos al modelo posterior V2.2.
\subsubsection{Executable semantics and the seven V1 property obligations}
\label{subsubsec:v1-executable-semantics}

The definitive \texttt{paper1-q3-v1} campaign evaluates the executable
\path{specs/tla/CrossShardCommit.tla} model, not the subsequent V2.2 model. Let $T$,
$S$, and $V$ denote the finite sets of transfers, shards, and validators.
The state vector is
\begin{equation}
\begin{split}
X=\langle &\mathit{status},\mathit{terminalStatus},\mathit{sourceShard},
\mathit{targetShard},\mathit{locked},\\
&\mathit{receiptOwner},\mathit{receiptUseCount},
\mathit{destinationCredit},\mathit{fundsReleased},\\
&\mathit{messages},\mathit{votes}\rangle.
\end{split}
\label{eq:editorial-v1-state}
\end{equation}

\texttt{Init} sets each transfer's status to \texttt{Pending}, terminal
status to \texttt{None}, lock/credit/release flags to false, receipt
uses to zero, and message/vote sets to empty. It chooses distinct source
and destination shards for each transfer. The executable transition
relation \texttt{Next} comprises \texttt{LockTransfer},
\texttt{ReleaseDelayedReceipt}, \texttt{ConsumeReceipt}, \texttt{CastVote},
\texttt{CommitTransfer}, \texttt{TimeoutTransfer}, and \texttt{Stutter},
with \texttt{CreditBeforeReceipt} and \texttt{CommitAfterAbort} guarded
by mutation modes and therefore disabled in the correct model. TLC
checks \(\mathit{Spec}=\mathit{Init}\land\Box[\mathit{Next}]_X\)
under the finite configurations supplied with the artifact.

In the correct model, \texttt{ConsumeReceipt} requires a locked transfer,
a non-delayed receipt message, zero previous uses, and no prior credit.
The same action increments the receipt-use counter and sets the Boolean
credit flag; this model does not expose an independently fallible
cross-shard destination-credit operation. \texttt{CommitTransfer}
requires a prepared transfer, destination credit, and at least the
configured quorum of votes. An enabled \texttt{TimeoutTransfer} in the
locked state marks abort and release in the same abstract transition.
These are substantial modeling restrictions, not implementation-level
proofs of distributed atomic commit.

The seven invariants checked in the V1 campaign are indexed below to
keep their full mathematical definitions readable. The index-to-source
correspondence is exact:
\begin{center}
\small
\begin{tabular}{ll}
$P_1$ & \texttt{NoReceiptReplay} \\
$P_2$ & \texttt{DestinationCreditRequiresValidReceipt} \\
$P_3$ & \texttt{DecisionConsistency} \\
$P_4$ & \texttt{EventuallyReleasedAfterTimeout} \\
$P_5$ & \texttt{QuorumRequired} \\
$P_6$ & \texttt{NoValueLossAtTermination} \\
$P_7$ & \texttt{TerminalStateIrreversibility}
\end{tabular}
\end{center}

Using the actual state variables of the executable V1 model, the seven predicates are
\begin{align}
P_1 &\equiv \forall t\in T:\mathit{receiptUseCount}[t]\leq1,\notag\\
P_2 &\equiv \forall t\in T:
  \mathit{destinationCredit}[t]\Rightarrow\notag\\
  &\qquad\bigl(\mathit{locked}[t]\land
  \mathit{receiptOwner}[t]=\mathit{targetShard}[t]\land\notag\\
  &\qquad\mathit{receiptUseCount}[t]=1\bigr),\notag\\
P_3 &\equiv \forall t\in T:
  \mathit{status}[t]=\texttt{Committed}\Rightarrow\notag\\
  &\qquad\neg\mathit{fundsReleased}[t],\notag\\
P_4 &\equiv \forall t\in T:
  \mathit{status}[t]=\texttt{Aborted}\Rightarrow\notag\\
  &\qquad\mathit{fundsReleased}[t],\notag\\
P_5 &\equiv \forall t\in T:
  \mathit{status}[t]=\texttt{Committed}\Rightarrow\notag\\
  &\qquad\lvert\mathit{votes}[t]\rvert\geq\mathit{Quorum},\notag\\
P_6 &\equiv \forall t\in T:\bigl(
  \mathit{status}[t]\in\{\texttt{Pending},\texttt{Locked},\texttt{Prepared}\}\notag\\
  &\qquad\lor (\mathit{status}[t]=\texttt{Committed}\land
     \mathit{destinationCredit}[t])\notag\\
  &\qquad\lor (\mathit{status}[t]=\texttt{Aborted}\land
     \mathit{fundsReleased}[t])\bigr),\notag\\
P_7 &\equiv \forall t\in T:\bigl(
     \mathit{terminalStatus}[t]=\texttt{None}\land\notag\\
  &\qquad\mathit{status}[t]\in
     \{\texttt{Pending},\texttt{Locked},\texttt{Prepared}\}\notag\\
  &\qquad\lor(\mathit{terminalStatus}[t]=\texttt{Committed}\land\notag\\
  &\qquad\mathit{status}[t]=\texttt{Committed})\notag\\
  &\qquad\lor(\mathit{terminalStatus}[t]=\texttt{Aborted}\land\notag\\
  &\qquad\mathit{status}[t]=\texttt{Aborted})\bigr).
\label{eq:editorial-v1-invariants}
\end{align}

These formulas use the actual state variables of the V1 TLA+ module. In
particular, $P_4$ is an invariant about the \texttt{Aborted} state rather
than a temporal eventuality. The formulas neither introduce new
verification outcomes nor retroactively alter the 1,272-task protocol.

Moreover,
several properties follow closely from action guards and simultaneous
state updates. Detecting a mutant that removes exactly such a guard
therefore does not independently establish implementation correctness,
non-vacuity of every property, or adequacy against broader defect classes.
The V1 state has a transfer-level terminal decision and Boolean
credit/release indicators; it cannot express two independently
irreversible decisions, one per shard, or failures between all concrete
ledger updates.

The later TLA+ V2.2 specification introduces distinct
\texttt{srcState}, \texttt{dstState}, and \texttt{decision} variables
and explicit \texttt{ApplySrc}/\texttt{ApplyDst} transitions. G2-P3
recorded three Java failure-injection traces with partial source debit
followed by an exact rollback, while V2.2 does not reverse an applied
commit debit. Thus neither one-to-one internal correspondence nor a
general Java--V2.2 refinement proof is claimed. The V2.2 source and
configurations support the semantic comparison with the Java witnesses.
The available V2.2 source copy lacks the original logs for the 52-case
matrix, so this manuscript reports no execution totals or verdicts from
that matrix. The Java G1--G2 observations are supplementary to the frozen
1,272-task V1 campaign.
